\label{sec:overview}

The AgriHome Vision Console is designed using a \textbf{layered architecture}. Each layer has a distinct role; layers interact through well-defined interfaces (HTTP APIs between client and server, device-key APIs to edge agents, internal service calls within the server, SQL and file I/O toward stores, optional outbound calls to inference services, and optional HTTP still pulls to LAN streamers). This separation improves modularity, testability, and the ability to scale or replace one concern (for example, swapping a vision provider or webcam path) without rewriting unrelated areas.

The system comprises \textbf{five primary layers}:

\begin{enumerate}
  \item \textbf{Presentation Layer} — user-facing experience in the browser or PWA: landing and marketing entry, authentication, dashboard, tray and plant views, mesh management, schedules, \textbf{Devices} navigation, tray Raspberry~Pi panel with Take Picture, charts, and uploads.
  \item \textbf{Application Layer} — server-side control: HTTP and GraphQL entry points (including \texttt{/api/raspberry-pi/*} and \texttt{/api/devices/*}), session and device-key verification, validation, routing of work to domain services, coordination of reads and writes to the data layer, invocation of the processing layer when needed, and optional HTTP streamer still fetches.
  \item \textbf{Processing Layer} — computer-vision and analytics \emph{adapters}, plus edge post-ingest hooks and the capture schedule runner: optional services that accept images or features and return structured labels, detections, or scores; schedule runners that enqueue \texttt{capture\_now} work for \texttt{raspberry-pi-edge} destinations.
  \item \textbf{Data Layer} — durable storage: relational database for domain and edge entities (\texttt{edge\_devices}, \texttt{edge\_device\_commands}, \texttt{capture\_pose\_sequences}, \texttt{capture\_poses}, \texttt{capture\_schedules}, \texttt{camera\_captures}, \texttt{plants}, \texttt{prediction\_results}, trays, meshes, reports), and object or file storage for imagery; optional vector index for similarity search when enabled.
  \item \textbf{Hardware Layer} — Raspberry~Pi hosts running Agri-Home/klipper (\texttt{fswebcam} / \texttt{camera-macros/save\_image.sh}), optional LAN HTTP streamer bases, and the \texttt{agrihome\_agent} bridge that registers, heartbeats, captures, and ingests.
\end{enumerate}

\subsection{Presentation layer description}

The Presentation Layer is the primary interface between \textbf{people} and the product. It implements responsive layouts for desktop and mobile, supports installable PWA behavior for field-adjacent use, and communicates with the Application Layer exclusively through HTTPS (including authenticated cookie-based sessions after sign-in).

Users register or sign in, navigate trays and plants, upload photos, trigger or view vision-assisted results, inspect monitoring history, configure schedules, manage meshes, and operate \textbf{edge devices} (list devices, edit Klipper/streamer URL, Take Picture from the tray Raspberry~Pi panel). The layer is responsible for presenting SRS-mandated capabilities such as health status, historical trends, dashboard summaries, and settings-oriented flows at a user-experience level; it does not embed direct database access.

\textbf{Key responsibilities:}
\begin{itemize}
  \item Capture user input (forms, file selection, navigation, Take Picture).
  \item Render charts, tables, imagery, device status, and structured errors clearly.
  \item Invoke application APIs and display responses or structured errors.
  \item Encapsulate client-only concerns (progressive enhancement, local caching of shell assets where applicable).
\end{itemize}

\subsection{Application layer description}

The Application Layer is the \textbf{central coordinator}. It receives every authenticated request from the Presentation Layer, interprets it, enforces authorization tied to the authenticated identity, validates payloads, delegates to domain-oriented modules (tray and plant lifecycle, mesh topology, schedules, monitoring logs, file references, edge device actions), and returns responses.

It also terminates \textbf{device-key} traffic from Raspberry~Pi agents (\texttt{/api/raspberry-pi/register}, heartbeat, ingest, commands, poses, capture-plan) and operator device actions (\texttt{/api/devices}). For Take Picture, this layer primarily queues \texttt{capture\_now} for the agent (\texttt{fswebcam}); when \texttt{klipper\_url} is set it may optionally fetch an HTTP streamer still first. It bridges to the Processing Layer for species/disease classification and schedule-driven capture enqueue.

\textbf{Key responsibilities:}
\begin{itemize}
  \item Expose a stable API surface for the console, edge agents, and future integrators.
  \item Maintain session integrity, ownership rules, and hashed device API keys.
  \item Orchestrate transactions across relational updates and file metadata (including multipart ingest).
  \item Isolate failures in optional backends so core monitoring remains usable when possible.
\end{itemize}

\subsection{Processing layer description}

The Processing Layer performs \textbf{compute-intensive or model-based} work and \textbf{asynchronous capture orchestration} that should not block the interactive UI thread. In AgriHome this includes, when configured, \textbf{species or leaf-level classification}, \textbf{tray-level vision}, the \textbf{edge disease hook} after Pi ingest, and the \textbf{capture schedule runner} that enqueues edge capture commands for destinations marked \texttt{raspberry-pi-edge}.

When no external model is available, the Application Layer may apply \textbf{deterministic or simulated} behavior for development and demonstration; the architecture still treats that path as ``processing'' logically because it produces structured inference output distinct from raw CRUD.

\textbf{Key responsibilities:}
\begin{itemize}
  \item Accept image or feature input in a contract agreed with the application tier.
  \item Run detection, embedding, or classification pipelines as implemented (including post-ingest disease analysis).
  \item Return normalized results (labels, confidence, geometry) to the Application Layer.
  \item Enqueue scheduled edge captures without requiring an operator to be present.
  \item Remain stateless with respect to long-term domain truth; persistence is always via the application and data layers.
\end{itemize}

\subsection{Data layer description}

The Data Layer is the \textbf{authoritative foundation} for AgriHome. The relational store holds entities and relationships that satisfy SRS concepts: plant health records over time, capture references, environmental or monitoring events as modeled, mesh membership, schedules, and Summer~2026 \textbf{edge} tables (\texttt{edge\_devices}, \texttt{edge\_device\_commands}, pose sequences/poses). Object storage holds binary assets (camera frames, user uploads) referenced by identifiers or URLs in the relational model.

Optional components such as a \textbf{vector database} participate in this layer when similarity or embedding search is required; they remain secondary to the relational system of record.

\textbf{Key responsibilities:}
\begin{itemize}
  \item Persist and retrieve structured domain and edge data with consistency expectations documented per operation.
  \item Store and stream large binaries efficiently.
  \item Enforce integrity through schema, constraints, and application-side validation.
  \item Support backup, migration, and operational monitoring consistent with deployment policy.
\end{itemize}

\subsection{Hardware layer description}

The Hardware Layer is the \textbf{physical capture plane}. Raspberry~Pi devices run Agri-Home/klipper with \texttt{fswebcam} and the \texttt{agrihome\_agent}, which auto-provisions into AgriHome, reports online status, and delivers imagery by answering queued commands (primary) or by exposing an optional HTTP streamer that the Application Layer can pull when the Pi is LAN-reachable from the server.

\textbf{Key responsibilities:}
\begin{itemize}
  \item Provide durable device identity and local agent configuration.
  \item Expose JPEG snapshot endpoints for server-side and agent-local capture.
  \item Execute pose-aware multi-angle capture when schedules or operators request it.
  \item Avoid autonomous plant-care actuation; motion fields may be stubbed pending hardware bring-up.
\end{itemize}

\subsection{Layered architecture figure}

Figure~\ref{fig:layer-arch} summarizes the five layers and primary external actors. Companion diagrams in Section~3 refine internal subsystems and data flow.

\begin{figure}[H]
\centering
\begin{tikzpicture}[
  font=\small,
  layer/.style={draw, rounded corners, minimum width=11.4cm, minimum height=1.05cm, align=left, inner sep=9pt},
  arr/.style={-{Stealth}, thick},
]
\node[layer, fill=leafgreen!18] (L1) {\textbf{Presentation Layer}\quad\textit{(browser / PWA; Devices; tray Pi panel)}};
\node[layer, below=0.38cm of L1, fill=blue!8] (L2) {\textbf{Application Layer}\quad\textit{(APIs, sessions, /api/raspberry-pi/*, /api/devices/*)}};
\node[layer, below=0.38cm of L2, fill=orange!14] (L3) {\textbf{Processing Layer}\quad\textit{(CV adapters; edge-disease-hook; schedule runner)}};
\node[layer, below=0.38cm of L3, fill=gray!14] (L4) {\textbf{Data Layer}\quad\textit{(PostgreSQL domain + edge tables; object storage)}};
\node[layer, below=0.38cm of L4, fill=teal!16] (L5) {\textbf{Hardware Layer}\quad\textit{(Pi; Klipper/fswebcam; agrihome\_agent)}};

\node[draw, circle, minimum size=0.6cm, left=0.55cm of L1, font=\scriptsize, inner sep=0pt] (u) {U};
\node[draw, rounded corners, font=\scriptsize, right=0.35cm of L1, align=center, inner sep=4pt] (idp) {Identity\\provider};

\draw[arr] (u) -- node[above, font=\scriptsize] {uses} (L1);
\draw[arr, dashed] (L1.east) -- (idp.west);
\draw[arr] (L1.south) -- (L2.north);
\draw[arr, dashed] (L2.south) -- node[right, font=\scriptsize] {when needed} (L3.north);
\draw[arr] (L3.south) -- (L4.north);
\draw[arr, bend left=18] (L5.west) to node[left, font=\scriptsize, align=center] {register\\heartbeat\\ingest} (L2.west);
\draw[arr, dashed, bend right=12] (L2.east) to node[right, font=\scriptsize] {snapshot} (L5.east);
\end{tikzpicture}
\caption{AgriHome Vision Console: five-layer architecture with Pi/Klipper edge trust boundaries.}
\label{fig:layer-arch}
\end{figure}

\subsection{System workflow overview}

A typical end-to-end path begins when a signed-in user opens a \textbf{tray} or \textbf{Devices} view in the Presentation Layer. The client requests data from the Application Layer, which reads the \textbf{Data Layer} (relational and, when needed, media URLs). For \textbf{Take Picture}, the Application Layer primarily queues a command that \texttt{agrihome\_agent} claims on heartbeat, captures via \texttt{fswebcam}/\texttt{save\_image.sh}, and multipart-ingests; if an optional streamer URL is reachable it may pull a still first. When vision-assisted analysis is configured, the Processing Layer's edge disease hook may classify the capture; results return as structured data and are written to \texttt{prediction\_results} / reports. Scheduled \texttt{raspberry-pi-edge} policies are advanced by the capture schedule runner. Updated state flows back to the Presentation Layer for display.
